\section{The translation class and the double commutator}\label{sec:commutator}

The preceding sections give the two filtration bounds for the same operator
$f'$.  We now compute its double commutator with $e$.  The calculation has two
parts: an elementary modification differentiates the universal degree-two
class into a difference of flag classes, and flag integration converts the
square of that difference back into $e$.

\subsection{The characteristic class under translation}

The universal projective bundle on $C\times\mathcal Y$ defines a universal
endomorphism bundle.  Let
\[
 e_{\mathcal Y}
   =\frac1{2n}\int_C\operatorname{ch}_2(\operatorname{End}E)
   \in H^2(\mathcal Y,\mathbb Q).
\]
Individual ordered flag lines $Q_i$ can be taken on local universal-bundle
charts; their ratios $Q_i\otimes Q_j^{-1}$ descend to the coarse moduli
space.  We use the global classes
\[
 y=c_1(Q_1\otimes Q_2^{-1}),\qquad
 t_C=x^*c_1(T_C),
\]
where the second is pulled back from $S$.  We retain the same notation for
the restrictions of these classes to a fiber $Y_{\lambda,u}$.

\begin{lemma}\label{calc:translation-formula}
For the translation $T:Y_{0,v}\to Y_{0,v}$ associated with
$\ell=(1,-1,0,\ldots,0)$ and every integer $k\ge0$,
\begin{equation}\label{calc:translation-polynomial}
 T^{k*}e_{\mathcal Y}=e_{\mathcal Y}+ky+k^2t_C.
\end{equation}
Consequently the logarithm from Section~\ref{sec:translation} satisfies
\begin{equation}\label{calc:logarithm}
 N(e_{\mathcal Y})=y
\end{equation}
on $H^*(\mathcal Y,\mathbb Q)$ and, by restriction, on every fiber.
\end{lemma}

\begin{proof}
Work first on $Y_{0,v}$.  Denote by $i:\Delta\hookrightarrow C\times Y_{0,v}$
the graph of the moving point.  Its normal line is the pullback of $x^*T_C$,
so
\[
 \mathcal O(\Delta)|_\Delta\simeq x^*T_C,
 \qquad c_1(N_\Delta)=t_C.
\]
The residue eigenvalues $v_i$ are distinct.  Consequently the Higgs operator
has uniquely labeled eigenline summands on the formal neighborhood of
$\Delta$.  This follows by lifting its simple residual eigenvalues and
their idempotent projectors to every finite order.  Multiplying a local
frame of the logarithmic canonical line by a unit does not change those
projectors.  Thus the formal summands are intrinsic, and any finite number
of their jets suffices for an elementary modification.

The $k$-fold translation replaces the $i$th formal line lattice by its twist
by $k\ell_i\Delta$.  On the $(i,j)$ block of $\operatorname{End}E$, the
lattice shift is therefore
\[
 m_{ij}=k(\ell_i-\ell_j).
\]
If a line block has restriction $F$ to $\Delta$, its K-theory change under
an integer shift $m$ is
\begin{equation}\label{calc:signed-jets}
 [\text{new block}]-[\text{old block}]
   =\sum_{a=1}^{m}[i_*(F\otimes N_\Delta^{\otimes a})].
\end{equation}
For $m<0$, the right side means the negative of the sum over $m<a\le0$;
for $m=0$ it is zero.  For positive $m$, the formula follows by successive
quotients of the lattices at pole orders $1,\ldots,m$.  For negative $m$,
apply the same quotient sequence between the smaller lattice and the
original one.  In particular $m=-1$ contributes $-[i_*F]$, as it must for
the kernel elementary transformation.

The resulting identity for $\operatorname{End}E$ is global in K-theory,
before applying a characteristic class.  Indeed the intrinsic formal
idempotents lie in the descended endomorphism algebra, so its $(i,j)$
blocks are global formal line bundles.  Identify the original and modified
endomorphism bundles away from $\Delta$, and choose a common sufficiently
small lattice inside both.  Their quotients by this lattice are supported
on a finite infinitesimal neighborhood of $\Delta$.  The formal block
projectors and the filtration by pole order give global finite
filtrations of these quotients.  Their associated graded terms are
graph pushforwards of the descended line bundles
$Q_i\otimes Q_j^{-1}\otimes N_\Delta^{\otimes a}$.  Subtracting the two
quotient classes cancels the common pole orders and leaves exactly the
signed sums in \eqref{calc:signed-jets}.  Additivity for these filtrations
therefore gives the identity in the Grothendieck group of coherent sheaves
on $C\times Y_{0,v}$; the common lattice also cancels.

Grothendieck--Riemann--Roch for the graph \cite[Section~15.2]{Fulton1998} gives
\[
 \operatorname{ch}(i_*F)
   =i_*\bigl(\operatorname{ch}(F)\operatorname{td}(N_\Delta)^{-1}\bigr).
\]
Hence the contribution of a summand in \eqref{calc:signed-jets} to
$\int_C\operatorname{ch}_2$ is
\[
 c_1(F)+(a-\tfrac12)t_C.
\]
This also exhibits the sign of the normal-bundle correction: the normal
line is the tangent line, and its inverse Todd class starts with
$1-t_C/2$.
For the signed sums just specified, the elementary identities
\[
 \sum_{a=1}^{m}1=m,
 \qquad
 \sum_{a=1}^{m}(a-\tfrac12)=\frac{m^2}{2}
\]
hold for every integer $m$.

Write $w_i=c_1(Q_i)$ on a universal-bundle chart.  Only $w_i-w_j$ and
expressions invariant under common translation of the $w_i$ will occur.
Summing the contributions for $F=Q_i\otimes Q_j^{-1}$ gives
\[
 T^{k*}e_{\mathcal Y}-e_{\mathcal Y}
 =\frac1{2n}\sum_{i,j}
   \left(k(\ell_i-\ell_j)(w_i-w_j)
       +\frac{k^2}{2}(\ell_i-\ell_j)^2t_C\right).
\]
Since $\sum_i\ell_i=0$,
\[
 \sum_{i,j}(\ell_i-\ell_j)(w_i-w_j)=2n\sum_i\ell_iw_i,
 \qquad
 \sum_{i,j}(\ell_i-\ell_j)^2=2n\sum_i\ell_i^2.
\]
For the chosen cocharacter these are $2ny$ and $4n$, respectively, proving
\eqref{calc:translation-polynomial}.  Since the lattice identity was global
on the coarse moduli space before applying Grothendieck--Riemann--Roch,
this is a global cohomological identity there.

Let $M$ be the integer used to define $N$.  Replace $k$ by $Mq$ in
\eqref{calc:translation-polynomial}.  Since $\mathsf T^{Mq}=\exp(MqN)$,
both sides are polynomial functions of the nonnegative integer $q$ with
values in a finite-dimensional cohomology group.  They agree as
polynomials.  Their linear coefficients are respectively $MN(e_{\mathcal
Y})$ and $My$, giving \eqref{calc:logarithm} on $Y_{0,v}$.
Corollary~\ref{log:restriction} transports this identity to
$H^*(\mathcal Y)$ and then to all its fibers: the three classes in
\eqref{calc:translation-polynomial} are restrictions of the global classes
just defined.
\end{proof}

The quadratic term in \eqref{calc:translation-polynomial} is consistent
with the behavior of the flag lines themselves.  Restricting a shifted
lattice to $\Delta$ gives
\[
 T^{k*}y=y+2kt_C,
 \qquad N(y)=2t_C,
 \qquad N(t_C)=0.
\]
In particular one must not assume $N(y)=0$ in the commutator calculation.

\subsection{The flag integral}

We record explicitly the permutation identity that evaluates the flag
correspondence.  It applies to projective bundles because every polynomial
used below is invariant under simultaneous translation of the roots.

\begin{lemma}\label{calc:flag-sum}
Let $\pi:\operatorname{Fl}(V)\to B$ be the full-flag bundle of a rank-$n$
complex vector bundle, and let $w_1,\ldots,w_n$ be its formal Chern roots.
For a polynomial $F$ in the first Chern classes of the ordered flag lines,
\[
 \pi_*\bigl(F(Q_\bullet)c_b(T_\pi)\bigr)
   =\sum_{\sigma\in\mathfrak S_n}
      F(w_{\sigma(1)},\ldots,w_{\sigma(n)}).
\]
If $\sum_i\ell_i=0$, then
\begin{equation}\label{calc:permutation}
 \sum_{\sigma\in\mathfrak S_n}
       \left(\sum_i\ell_iw_{\sigma(i)}\right)^2
 =n(n-2)!\|\ell\|^2
    \left(\sum_iw_i^2-\frac{(\sum_iw_i)^2}{n}\right).
\end{equation}
Both identities, when their expressions are invariant under common
translation of the roots, hold for the associated flag bundle of a
principal $\mathrm{PGL}_n$-bundle with rational coefficients.
\end{lemma}

\begin{proof}
The projective-bundle pushforward formula is a sum over the roots, with
denominator the product of the other-root differences.  Iterating it along
the full-flag tower gives the sum over ordered roots, with denominator the
tangent Euler class at that ordering.  Multiplication by $c_b(T_\pi)$
cancels this denominator, giving the first identity.  One may perform this
calculation after passing to a splitting space and inverting root
differences; the resulting equality is a polynomial identity in the
universal Chern classes, so it holds without the inversion.  The
projective-bundle pushforward is characterized by its Chern-polynomial
relation and the normalization that the top power of the hyperplane class
integrates to one; using the complex tangent Euler class fixes the
orientation signs in the displayed cancellation.

For \eqref{calc:permutation}, a fixed root occurs $(n-1)!$ times at a fixed
position, and a fixed ordered pair of distinct roots occurs $(n-2)!$ times
at two prescribed distinct positions.  Since
$\sum_{i\ne j}\ell_i\ell_j=-\|\ell\|^2$, expansion of the square gives
\[
 \begin{split}
 &(n-1)!\|\ell\|^2\sum_iw_i^2
 -(n-2)!\|\ell\|^2\sum_{i\ne j}w_iw_j\\
 &\hspace{15mm}
 =n(n-2)!\|\ell\|^2
    \left(\sum_iw_i^2-\frac{(\sum_iw_i)^2}{n}\right).
 \end{split}
\]
Finally, rational pullback from $B\mathrm{PGL}_n$ to $B\mathrm{GL}_n$ is
injective: on a maximal torus it is the inclusion of symmetric polynomials
in the root differences into symmetric polynomials in all roots.
Universal characteristic-class identities invariant under common root
translation can therefore be checked after this pullback, proving the
last assertion.
\end{proof}

\subsection{The common lowering operator}

\begin{proposition}\label{calc:double-commutator}
Let $\delta=\deg(x:S\to C)$.  The operator $f'$ of
\eqref{trans:fprime} satisfies
\begin{equation}\label{calc:scalar}
 [[f',e],e]=8(-1)^b n(n-2)!\delta\,e.
\end{equation}
Consequently
\begin{equation}\label{calc:f}
 f=\frac{(-1)^{b+1}}{4n(n-2)!\delta}\,f'
\end{equation}
has cohomological degree $-2$, lowers perversity by two and Betti weights
by four, and satisfies $[[f,e],e]=-2e$.
\end{proposition}

\begin{proof}
Work on the scalar correspondence at $\lambda=0$.  The restriction of
$e_{\mathcal Y}$ along $\iota:Z_0\hookrightarrow Y_{0,0}$ equals
$\pi^*\operatorname{pr}_0^*e$: the universal projective characteristic
class is unchanged by the line twist $J$ in the scalar product
identification.  Projection formulas therefore show that the maps outside
$N_{0,0}^2$ in $f'$ intertwine multiplication by the corresponding classes
$e$ and $e_{\mathcal Y}$.

For a degree-zero derivation $N$ and an even class $a$, write $m_a$ for
multiplication by $a$.  Then $[N,m_a]=m_{N(a)}$, and direct expansion gives
\[
 [[N^2,m_a],m_a]=2m_{N(a)^2}.
\]
Indeed the first commutator is
$m_{N^2(a)}+2m_{N(a)}N$; the first term commutes with $m_a$, while the
second contributes $2m_{N(a)^2}$.  Thus no vanishing assertion about
$N^2(a)$ is needed.  By Lemma~\ref{calc:translation-formula},
\begin{equation}\label{calc:middle-commutator}
 [[f',e],e]
 =2\operatorname{pr}_{0,*}\mathsf A_0
       (y^2\cup-)\mathsf B_0\operatorname{pr}_0^*.
\end{equation}

The normal bundle of $\iota$ is $T_\pi^*$.  Self-intersection therefore
replaces $\iota^*\iota_*$ in \eqref{calc:middle-commutator} by
\[
 c_b(T_\pi^*)\cup-=(-1)^b c_b(T_\pi)\cup-.
\]
For a class pulled back from $X$, the projection formula and
Lemma~\ref{calc:flag-sum} reduce the flag integral to the symmetric class
\[
 \begin{split}
 \sum_{\sigma\in\mathfrak S_n}
     (w_{\sigma(1)}-w_{\sigma(2)})^2
 &=2n(n-2)!
    \left(\sum_iw_i^2-\frac{(\sum_iw_i)^2}{n}\right)\\
 &=4n(n-2)!\,\widetilde{\operatorname{ch}}_2(E_{x(s)}).
 \end{split}
\]
Here $E_{x(s)}$ denotes the universal projective data restricted to the
graph of $x$; the normalized expression is well defined without an honest
universal vector bundle.

Under $O_0\simeq S\times X$, this characteristic class is the pullback of
$\widetilde{\operatorname{ch}}_2(E)$ from $C\times X$ along
$x\times\operatorname{id}_X$.  Tensoring by $J$ contributes nothing to it.
Integration along $S$ multiplies integration along $C$ by $\delta$:
\[
 \operatorname{pr}_{0,*}
      \widetilde{\operatorname{ch}}_2(E_{x(s)})=\delta e.
\]
This follows directly from the degree formula $x_*[S]=\delta[C]$, or from
the $C$-degree-two term in the K\"unneth decomposition.  The other curve
degrees integrate to zero.
Combining this equality, the factor $2$ in
\eqref{calc:middle-commutator}, and the self-intersection sign proves
\eqref{calc:scalar}.

The integer $\delta$ is positive because $x$ is a nonempty finite
\'{e}tale cover of connected projective curves.  Thus
\eqref{calc:f} is a nonzero rational rescaling and gives the stated double
commutator.  The degree of $f'$ is $-2$, since the two flag shifts cancel
and integration over $S$ has degree $-2$.  The filtration bounds are
Propositions~\ref{spec:perverse-lowering} and~\ref{trans:weights}; they are unchanged by the rescaling.
\end{proof}
